% Page 101

\seite{101}

ausfallen mußte, wurde sie mit der Hauptver\ohb
sammlung verbunden. Anschließend wurde die\ob
Tagesordnung erledigt, die recht ihre Vielseitig keit\ob
über drei Stunden beanspruchte. Um 22 Uhr konnte\ob
der Vorsitzende die anregend verlaufende Versamm\ohb
lung schließen. Von einer Familienfeier in der\ob
\unsicher Mutter wurde Abstand genommen. Am\ob
1.\ März wird eine alljährliche Heldengedächtnis\ohb
feier stattfinden.

\margmark{20. Februar}%
Ihren 33.\ Geburtstag feiert heute die älteste\ob
Einwohnerin unserer Gemeinde, die Witwe Magda\ohb
lena Peters. Die Jubilarin, die dieser Tage in voller\ob
körperlicher und geistiger Frische das große Dauer\ohb
fest am Morgen-Lebensabend feiert und die Strapazen\ob
genügt zurückzublicken. Sie als älteste Bürgerin, wie Dokt\ohb
\unsicher der damaligen Münzen, hat die drei Kriege\ob
1864/66, 1870/71 und die Kriegsjahre. Im Kriege 1870/71 miterlebt und mit\ohb
gemacht. Ihren einzigen Sohn hat sie im Welt\ohb
kriege dem Vaterlande auf dem Felde der Ehre ge\ohb
opfert. Möge es ihr vergönnt sein, noch viele Jahre\ob
bester Gesundheit entgegenzukommen!

\margmark{5. März}%
In den letzten Wochen wurden in der Jagdgemein\ohb
schaft unserer Pfarrei fünfzehn Wildschweine ge\ohb
schossen. Es war bisher zu erstmaligen Vorkommen der\ob
dabei gelang es Lehrer Müller von Schleid auf der\ob
Heilenbacher Jagd einen Keiler von ca. 2 Zentnern\ob
zu erlegen. Der Sefferweicher Jagdpächter\unsicher{} Darsbach\ob
\unsicher von Leimischof im Rhein von 70 Pfund\ob
und ein Elf\unsicher Hauer erlegte Adolf\unsicher{} Par\ob
von Neidenbach einen Keiler von 170 Pfund.

\margmark{12. März}%
Generalversammlung\unsicher{} 1 Uhr bei\unsicher{} der Gebrüder\unsicher\ob
Müller in Seffern\unsicher{} auf das\unsicher{}\unsicher des\unsicher\ob
groß\unsicher{} ein Vorstoß\unsicher abzuhalten\unsicher
\pend
